\chapter{Social Epistemology, Standpoints, and Values}
\label{ch:social}

\begin{chapterguide}
This chapter asks whether social location and institutional organization
undermine or improve objectivity. It treats science as a collective achievement,
examines feminist standpoint theory and situated knowledge, and distinguishes
values that select questions from values that distort evidence. The proposed
ideal is public, structured criticism rather than a view from nowhere.
\end{chapterguide}

No individual scientist can reproduce all the measurements, derivations,
computer codes, field observations, and historical records on which a mature
science depends. Scientific knowledge is therefore social in an unavoidable
sense. The question is whether the social character is a source of objectivity,
a source of bias, or both.

The answer is both. Division of labour allows specialization and independent
checking. It also creates authority gradients and dependencies that can hide
errors. Peer review can expose mistakes. It can also enforce conformity.
Funding can support long-term inquiry. It can also redirect attention toward
profitable or politically convenient questions. Social structure is not an
external background to epistemology. It is part of the machinery by which
evidence is produced and evaluated.

\section{The social organization of knowledge}

Social epistemology studies how people pursue truth with and through others.
Alvin Goldman emphasizes testimony, expertise, and the conditions under which
social processes are reliable \citep{goldman1999}. A person often knows
something because another person measured it, and the trustworthiness of that
testimony depends on competence, incentives, access to criticism, and a record
of past performance.

The division of cognitive labour can improve objectivity in at least four ways:

\begin{enumerate}
\item specialists develop techniques too complex for one investigator;
\item independent groups can detect local mistakes;
\item different backgrounds make different anomalies salient;
\item public records allow later investigators to audit procedures.
\end{enumerate}

But social dependence also creates vulnerabilities. A community can coordinate
around a false assumption. A prestigious laboratory can make contrary evidence
hard to publish. Researchers may replicate a result only when it is fashionable.
A field can measure what is easy rather than what is important.

\RCR\ therefore evaluates institutions by their \emph{correction capacity}:
the ability to expose a claim to competent, motivated, and sufficiently
independent criticism. Consensus is evidence only when the process generating
the consensus has a reasonable relation to truth.

\section{Merton's norms and their limits}

Robert Merton described norms associated with the scientific ethos:
communal sharing of knowledge, universalistic assessment, disinterestedness,
and organized skepticism \citep{merton1942}. The norms are ideals, not
descriptions of every laboratory. They express a valuable thought: scientific
claims should be judged by reasons that are not supposed to depend on the
speaker's identity or patronage.

The norms require supplementation. Universalistic assessment can become blind
to how research questions were selected. Communal sharing can conflict with
privacy or intellectual property. Disinterestedness does not mean that
scientists have no values; it means that personal or institutional interests
should not decide the evidential conclusion. Organized skepticism can be
distributed unevenly, with critics demanding impossible standards from
outsiders and weak standards from insiders.

\section{Longino and transformative criticism}

Helen Longino argues that objectivity is not achieved by eliminating social
values from inquiry. It is strengthened by critical interaction under
appropriate social conditions \citep{longino1990}. The relevant conditions
include recognized venues for criticism, uptake of criticism, public
standards, and a degree of equality in intellectual authority.

This account turns disagreement into an epistemic resource. A criticism is not
valuable merely because it is dissenting. It must identify a reason to doubt,
reveal a hidden assumption, offer an alternative, or show that a result does
not travel. The community must also be able to respond rather than ignore the
criticism.

\begin{principlebox}
\textbf{Critical interaction principle.} A scientific community is more
objective when relevant criticism can be expressed, understood, answered, and
incorporated without the critic having to possess the dominant group's
institutional status.
\end{principlebox}

The principle is compatible with expertise. Equality of intellectual authority
does not mean that every claim has equal evidential weight. It means that
access to criticism should not be determined solely by status, and that
expertise itself should be accountable to public reasons.

\section{Standpoint theory and situated knowledge}

Feminist standpoint theory argues that social positions can shape what is
visible and can sometimes provide epistemic advantages concerning relations of
power. Sandra Harding connects standpoint theory to the idea that inquiry
should begin from lives and activities that dominant institutions have treated
as peripheral \citep{harding1991}. Donna Haraway emphasizes situated
knowledges: every perspective is partial, and objectivity requires recognizing
the location and limits of the knower \citep{haraway1988}.

The strongest version is not that every member of an oppressed group
automatically knows the truth. Social position can generate insight, but it
can also generate partiality, pressure, or exclusion. Epistemic advantage is
an achievement of critical reflection and collective inquiry, not a biological
property.

Standpoint insights matter because a theory can fit the experience of a
dominant group while missing the mechanisms affecting others. Research on
workplace safety, medical treatment, domestic labour, disability, and
discrimination can improve when people who bear the consequences participate in
question formation and interpretation. This is not a concession that evidence
is subjective. It is an argument that the range of relevant evidence is
socially organized.

\section{Values in science}

The value-free ideal is ambiguous. There is a legitimate distinction between
values that determine what a result is and values that guide which questions
are worth asking. Scientific work requires choices about significance,
acceptable risk, measurement cost, privacy, classification, and the balance
between false positives and false negatives.

Heather Douglas argues that values can have a legitimate role in decisions
about acceptable error, especially where evidence informs policy
\citep{douglas2009}. Elizabeth Anderson shows how value judgments can guide
concept formation and interpretation without making the evidence itself a
matter of preference \citep{anderson2004}. Max Weber distinguished empirical
analysis from value relevance: values can select a problem without deciding
the factual answer \citep{weber1949}.

\begin{table}[ht]
\centering
\smalltable
\begin{tabularx}{\textwidth}{p{0.28\textwidth}YY}
\toprule
Role of value & Legitimate function & Epistemic risk\\
\midrule
Question selection & Decide what matters to investigate & Neglect of less visible populations\\
Concept formation & Choose distinctions relevant to a problem & Smuggling a conclusion into a definition\\
Loss function & Weigh false positives and false negatives & Treating a policy preference as a fact\\
Evidence interpretation & Decide when uncertainty warrants action & Confirmation bias and motivated reasoning\\
Institutional design & Set rules for access and correction & Exclusion, capture, and authority abuse\\
\bottomrule
\end{tabularx}
\caption{Values can enter different stages of inquiry. Their presence does not
make all stages subjective, but it makes transparency necessary.}
\label{tab:values}
\end{table}

The key distinction is between \emph{constitutive} and \emph{evidential}
values. Constitutive values help determine what to study, how to measure it,
and what risks to accept. Evidential values would make a claim count as true
because it serves an interest, regardless of the evidence. The first are
unavoidable and can be legitimate. The second corrupt objectivity.

\section{When institutions corrupt inquiry}

A research system becomes less objective when it systematically weakens the
routes by which errors can be found. Warning signs include:

\begin{itemize}
\item selective publication and undisclosed outcome switching;
\item opaque data processing and unavailable code;
\item conflicts of interest that affect design or interpretation;
\item exclusion of relevant populations from research;
\item status-based dismissal of criticism;
\item incentives that reward novelty while penalizing correction.
\end{itemize}

These are not merely ethical defects. They are epistemic defects because they
alter what evidence becomes visible and how it is weighted. A conflict of
interest does not prove a conclusion false, but it changes the expected
reliability of the process and should trigger stronger independent checks.

\section{Objectivity as social self-correction}

The view from nowhere is impossible. A more realistic ideal is a community that
can locate its own partiality. \RCR\ calls this \emph{reflexive objectivity}.
It has four components:

\begin{enumerate}
\item the perspective and values of inquiry are stated;
\item affected and relevantly situated people can challenge the framing;
\item independent criticism can inspect the evidence and analysis;
\item revisions are judged by their improved contact with the target, not only
by their political convenience.
\end{enumerate}

Reflexive objectivity is not guaranteed by diversity alone. A diverse group can
share the same measurement error. Nor is objectivity guaranteed by
disinterestedness alone. A narrow research agenda can be very carefully
pursued. The point is to connect social inclusion with world-directed tests.

\section*{Chapter summary}

Scientific knowledge is socially produced, and social organization can improve
or corrupt objectivity. Division of labour, expertise, criticism, and public
records create correction capacity. Merton's norms are valuable ideals but
require attention to power and question selection. Longino, Harding, and
Haraway show why critical pluralism and situated perspectives can reveal
ignored phenomena. Values legitimately shape questions, concepts, risk
thresholds, and institutions; they should not decide evidential conclusions by
fiat. Objectivity is therefore reflexive: it includes the ability of inquiry to
expose its own partiality.

\begin{discussion}
\begin{enumerate}
\item Is a scientific community objective if it is highly diverse but has no
effective correction mechanisms?
\item Which values should be public in a medical or policy study?
\item How can standpoint insights be tested without reducing them to identity
claims?
\end{enumerate}
\end{discussion}

\begin{furtherreading}
Read Longino (1990), Harding (1991), Haraway (1988), Anderson (2004),
Douglas (2009), Goldman (1999), and the literature on social dimensions of
scientific knowledge.
\end{furtherreading}

